\subsubsection{QAOA Ingredients}\label{sec:qaoa-ingredients}

In this section, we will see the high-level architecture of the Quantum Approximate Optimization Algorithm (QAOA). We will discuss \textbf{what components QAOA needs}, before explaining how each one works.

\highspace
In general, we can think of QAOA as a recipe. To make a dish, we need ingredients. To run QAOA, we need ingredients as well. The ingredients of QAOA are:
\begin{itemize}
    \item Optimization problem (the problem we want to solve)
    \item Cost Hamiltonian $H_C$
    \item Mixer Hamiltonian $H_M$
    \item Quantum circuit
    \item Classical optimizer
    \item Best parameters $(\gamma, \beta)$
    \item Approximate optimal solution (the solution we get from QAOA)
\end{itemize}
Let's go through each of these ingredients in detail.
\begin{enumerate}
    \item \important{Build the Cost Hamiltonian}. We convert our optimization problem into $H_C$. The ground state of the Cost Hamiltonian should be the solution of the problem we want to solve. In other words, it answers the question: ``\emph{what are we trying to optimize?}''.
    
    In optimization we usually define a \textbf{cost function} $C(x)$ that we want to minimize. QAOA replaces that classical cost function with $H_C$. So, instead of minimizing $C(x)$, we \textbf{minimize the expected energy} of $H_C$:
    \begin{equation*}
        \langle H_C \rangle = \bra{\psi} H_C \ket{\psi}
    \end{equation*}

    For example, for every possible solution:
    \begin{center}
        \begin{tabular}{@{} c c @{}}
            \toprule
            \textbf{Solution} & \textbf{Energy} \\
            \midrule
            A & 12 \\
            B & 7 \\
            C & \textbf{1} \\
            D & 9 \\
            \bottomrule
        \end{tabular}
    \end{center}
    The \textbf{lowest energy corresponds to the optimal solution}. So the Cost Hamiltonian simply answers: ``\emph{how good is this solution?}''.


    \item \important{Build the Mixer Hamiltonian}. Now suppose the quantum computer is currently representing solution B. \emph{How can it move toward solution C?} It \hl{needs a mechanism that changes the quantum state.} That mechanism is the \textbf{Mixer Hamiltonian}. It doesn't know what the best solution is, it just moves the quantum state around.
    
    In other words, the \definition{Mixer Hamiltonian} \textbf{explores the space of possible quantum states}. It tries to solve the local minima problem by moving the quantum state around.

    \textcolor{Green3}{\faIcon{question-circle} \textbf{Why both the Cost Hamiltonian and the Mixer Hamiltonian?}} Because they do different jobs. The Cost Hamiltonian \textbf{evaluates} how good a solution is, while the Mixer Hamiltonian \textbf{explores} the space of possible solutions. The two work together to find the best solution.
    \begin{itemize}
        \item Cost Hamiltonian: \emph{How good is this solution?}
        \item Mixer Hamiltonian: \emph{What other solutions should we try?}
    \end{itemize}

    
    \item \important{Construct the Quantum Circuit}. Once we have both Hamiltonians, they must be \textbf{transformed into an actual quantum circuit}. This raises an obvious question; if Hamiltonians are matrices and quantum computers execute gates, \emph{how do we convert one into the other?} The answer will be discussed in the next sections.
    
    In general, we will see that:
    \begin{equation*}
        U = e^{-i \alpha H}
    \end{equation*}
    Allows us to convert a Hamiltonian into a unitary operator, which can then be implemented as quantum gates. However, for now, just know that we can convert Hamiltonians into unitary operators, and then into quantum circuits.


    \item \important{Use a Classical Optimizer}. Just like VQE, QAOA is a \textbf{hybrid algorithm}. The quantum computer alone is \textbf{not enough}. So we need to choose a classical optimizer and begin to iterate. \emph{\textbf{Why?}} Because the circuit contains unknown parameters $(\gamma, \beta)$ that represent the angles of rotation for the Cost and Mixer Hamiltonians (we will see this in the next sections). The optimizer decides:
    \begin{itemize}
        \item Which values to try;
        \item Whether the new result is better;
        \item How to update the parameters.
    \end{itemize}
    

    \item \important{Optimization Objective}. Finally, \emph{\textbf{what are we minimizing?}} Exactly the same objective as VQE: the expected energy of the Cost Hamiltonian. The classical optimizer will try to find the best parameters $(\gamma, \beta)$ that minimize the expected energy of $H_C$:
    \begin{equation*}
        \langle H_C \rangle = \bra{\psi(\gamma, \beta)} H_C \ket{\psi(\gamma, \beta)}
        \quad
        \to
        \quad
        \argmin_{\theta} \bra{\psi(\theta)} H_C \ket{\psi(\theta)}
    \end{equation*}
    Where the parameters $\theta$ are the angles $(\gamma, \beta)$ of the Cost and Mixer Hamiltonians.

    Therefore, QAOA is searching for \hl{the values of $\gamma$ and $\beta$ that produce the quantum state with the lowest expected energy.} Since the Cost Hamiltonian was built so that the lowest energy corresponds to the optimal solution, this means that QAOA is searching for the best solution to the optimization problem.
\end{enumerate}
So to summarize, the entire algorithm can be summarized as follows:
\begin{enumerate}
    \item Optimization problem
    \item Build Cost Hamiltonian $H_C$
    \item Add Mixer Hamiltonian $H_M$
    \item Build parameterized circuit
    \item Choose $(\gamma, \beta)$
    \item Run circuit
    \item Measure
    \item Compute $\langle \psi | H_C | \psi \rangle$
    \item Classical optimizer
    \item Update $(\gamma, \beta)$
    \item Repeat
\end{enumerate}
Notice that this is essentially the \textbf{same hybrid optimization loop as VQE}. The only major difference is the \textbf{structure of the ansatz}: in VQE the ansatz is largely user-defined, whereas in QAOA the ansatz is \textbf{completely determined by the Cost Hamiltonian and the Mixer Hamiltonian}. That distinction will become the central topic of the next section.